\section{Comparison of \ours with State-of-the-art of Few-Shot Learning Models and Fine-Tuned Baselines}
\label{sec:sota}

In this section, we present a comprehensive comparison of the performance of \ours with other few-shot models and fine-tuned baselines across a range of tasks, including mathematical reasoning and programming tasks. Tables \ref{tab:pie_comparison} and \ref{tab:gsm-8k_comparison} display the performance of these models on the PIE dataset and GSM tasks, respectively. Our analysis demonstrates the effectiveness of different model architectures and training techniques in tackling complex problems.


\begin{table}[ht]
    \centering
    \begin{tabular}{llr}
        \toprule
        \multicolumn{2}{l}{\textbf{Method}} & \multicolumn{1}{c}{\textbf{Solve Rate}} \\
        \midrule
        \multirow{1}{*}{\citet{cobbe2021training}} & OpenAI 6B & 20.0 \\
        \multirow{1}{*}{\citet{chainofthought}} & CoT w/ \codex & 65.6 \\ \midrule
        \multirow{4}{*}{\citet{gao2022pal}} & PaL w/ \codex & 72.0 \\
   & PaL w/ GPT-3 & 52.0 \\
        & PaL w/ \gptt & 56.8 \\
        & PaL w/ \chatgpt & 74.2 \\
        & PaL w/ \gptf & 93.3 \\
        \multirow{2}{*}{\citet{Welleck2022SelfCorrect}} & Self-Correct w/ GPT-3 & 45.9 \\ 
         & Self-Correct (fine-tuned) & 24.3 \\ 
        \midrule
        \multirow{4}{*}{This work} & \ours w/ GPT-3 & \textbf{55.7} \\
        & \ours w/ \gptt & \textbf{62.4} \\
        & \ours w/ \chatgpt & \textbf{75.1} \\
        & \ours w/ \gptf & \textbf{94.5} \\
        \bottomrule
    \end{tabular}
\caption{Performance comparison of models on math reasoning (\gsm).}

    \label{tab:gsm-8k_comparison}
\end{table}


\begin{table}[ht]
    \centering

    \begin{tabular}{llr}
        \toprule
        \multicolumn{2}{l}{\textbf{Method}} & \multicolumn{1}{c}{\pctOpt)} \\
        \midrule
        \multirow{1}{*}{\citet{puri2021codenet}} & \textbf{Human References} & 38.2 \\ \midrule
      \multirow{4}{*}{OpenAI Models: \citet{openai_blogpost,openai2023gpt4}}  & \codex & 13.1 \\
        & \gptt & 14.8 \\
        & \chatgpt & 22.2 \\
        & \gptf & 27.3 \\
        \midrule
        \multirow{1}{*}{\citet{nijkamp2022codegen}} & \codegen-16B & 1.1 \\
        \midrule
        \multirow{3}{*}{\citet{scalene}} & \scalene & 1.4 \\
        & \scalene (\bestof{16}) & 12.6 \\
        & \scalene (\bestof{32}) & 19.6 \\
        \midrule
        \multirow{8}{*}{\citet{pie}} & \pie-2B & 4.4 \\
        & \pie-2B (\bestof{16}) & 21.1 \\
        & \pie-2B (\bestof{32}) & 26.3 \\
        & \pie-16B & 4.4 \\
        & \pie-16B (\bestof{16}) & 22.4 \\
        & \pie-16B (\bestof{32}) & 26.6 \\
       & \pie-Few-shot (\bestof{16}) & 35.2 \\
        & \pie-Few-shot (\bestof{32}) & \textbf{38.3} \\
        \midrule
       \multirow{3}{*}{This work} & \ours w/ \gptt & 23.0 \\
        & \ours w/ \chatgpt & 26.7 \\
        & \ours w/ \gptf & 36.0 \\
        \bottomrule
    \end{tabular}
    \caption{Performance comparison of various models on the PIE dataset in terms of the percentage of programs optimized (\pctOpt). The table includes human references, baseline models, fine-tuned \pie-2B and \pie-16B models, and our proposed model (\ours) using different LLMs. Notably, \ours achieves superior performance while using only 4 samples at most, significantly fewer than the 16 and 32 samples employed by other models. Scalene, an off-the-shelf optimizer, uses instruction tuning with Codex and serves as a comparison point.}
    \label{tab:pie_comparison}
\end{table}

